\chapter{The Probabilistic Method}

\begin{center}\itshape
To prove that a haystack contains a needle,\\
show that a straw picked at random is a needle with positive probability.
\end{center}

Some objects are easy to describe and hard to find: a colouring of a large complete graph with no big monochromatic
clique, an assignment satisfying a formula whose clauses barely overlap, a sparse graph with a giant connected
piece. Paul Erd\H os discovered in 1947 that one can prove such objects exist without constructing them: choose the
object at random, and show that it has the required property with positive probability. Since then the
\emph{probabilistic method} has become one of the main tools of combinatorics and theoretical computer science.

The chapter climbs in sophistication. The \emph{first moment method} uses only linearity of expectation: if the expected
number of bad configurations is less than $1$, some outcome has none. \emph{Alterations} first take a random object, then
repair it. The \emph{second moment method} shows that a random variable is positive with high probability, which
locates thresholds in random graphs. The \emph{Lov\'asz local lemma} handles many bad events that are each rare and
depend on few others, even when the union bound fails. Then we follow the random graph $G(n,\frac cn)$ through its phase
transition: for $c<1$ all components are small, for $c>1$ a giant one appears, and the branching processes of Chapter~3
explain why. Finally, \emph{conditional expectations} turn some of these existence proofs into fast deterministic
algorithms.

\begin{think}
Among any $6$ people there are $3$ mutual friends or $3$ mutual strangers. Is there a party of $1000$ people with no $20$
mutual friends and no $20$ mutual strangers? If so, could you describe the friendships explicitly?
\end{think}

\section{The first moment method}

Two facts of Chapter~2 do all the work: if $X$ is a random variable with values in $\N$ and $\mathbb E[X]<1$, then $P(X=0)>0$; and any random
variable takes, with positive probability, a value at least its mean. So an object with the desired property exists.

The \emph{Ramsey number} $R(k)$ is the smallest $n$ such that every colouring of the edges of the complete graph $K_n$ in red and blue
contains a monochromatic $K_k$ (a set of $k$ vertices whose $\binom k2$ edges all have the same colour).

\begin{proposition}\label{prop:ramseyupper}
$R(3)=6$, and for all $k\ge2$, $R(k)\le\binom{2k-2}{k-1}<4^{k-1}$.
\end{proposition}
\begin{proof}
Let $R(k,l)$ be the smallest $n$ such that every red--blue colouring of $K_n$ has a red $K_k$ or a blue $K_l$. We show $R(k,l)\le R(k-1,l)+R(k,l-1)$; with
$R(k,2)=k$ and $R(2,l)=l$, induction and Pascal's rule give $R(k,l)\le\binom{k+l-2}{k-1}$. Take $n=R(k-1,l)+R(k,l-1)$ and a vertex $v$ of $K_n$. Among the other
$n-1$ vertices, either at least $R(k-1,l)$ are joined to $v$ in red, or at least $R(k,l-1)$ in blue. In the first case, these vertices contain a blue $K_l$, or
a red $K_{k-1}$ which together with $v$ forms a red $K_k$; the second case is symmetric. For $k=3$ this gives $R(3)\le6$, and the colouring of $K_5$
of Figure~\ref{fig:ramsey} (the sides of a pentagon red, the diagonals blue) has no monochromatic triangle, so $R(3)=6$.
\end{proof}

\begin{figure}[ht]
\centering
\fig{ch12-ramsey}
\caption{Left: a colouring of $K_5$ without a monochromatic triangle. Right: a random colouring of $K_6$; it has a monochromatic
triangle (thick), as every colouring of $K_6$ must.}
\label{fig:ramsey}
\end{figure}

\begin{theorem}[Erd\H os, 1947]\label{thm:erdosramsey}
If $\binom nk2^{1-\binom k2}<1$, then $R(k)>n$. In particular $R(k)>2^{k/2}$ for $k\ge3$.
\end{theorem}
\begin{proof}
Colour each edge of $K_n$ red or blue independently with probability $\frac12$. A given set of $k$ vertices is monochromatic with probability
$2\cdot2^{-\binom k2}$, so the expected number of monochromatic $K_k$ is $\binom nk2^{1-\binom k2}<1$. Hence some colouring has none, and $R(k)>n$. For $n=\lfloor2^{k/2}\rfloor$,
\[
\begin{aligned}\binom nk2^{1-\binom k2}&\le\frac{n^k}{k!}2^{1-\frac{k(k-1)}2}\eqbr\le\frac{2^{k^2/2}}{k!}2^{1-\frac{k^2}2+\frac k2}=\frac{2^{1+k/2}}{k!},\end{aligned}
\]
which is $\frac{2^{2.5}}6<1$ for $k=3$ and decreases with $k$, since each step multiplies it by $\frac{\sqrt2}{k+1}<1$.
\end{proof}

So $\sqrt2^{\,k}<R(k)<4^k$. For $k=10$ the theorem gives $R(10)>100$, and for $k=20$, $R(20)>5817$. Nobody knows how to write down an explicit colouring
anywhere near these bounds: the best explicit constructions are far smaller for large $k$, and Erd\H os offered a prize for one of
size $(1+\varepsilon)^k$. Exact values are known only for $k\le4$ ($R(4)=18$); for $k=5$ it is known that $43\le R(5)\le46$. The upper bound $4^k$ was improved to
$(4-\varepsilon)^k$ only in 2023, by Campos, Griffiths, Morris and Sahasrabudhe. Tome~3, Appendix~B takes up Ramsey numbers from the combinatorial side:
the proof that $R(4)=18$, Tur\'an's theorem on dense graphs without cliques, and a probabilistic proof of it in the spirit of this chapter.

\begin{proposition}[large cuts]\label{prop:maxcut}
Every graph with $m$ edges has a \emph{cut} with at least $\frac m2$ edges: its vertices can be split into two sets so that at least $\frac m2$
edges join the two sets.
\end{proposition}
\begin{proof}
Put each vertex in one of the two sets by an independent fair coin. Each edge is cut with probability $\frac12$, so the expected
number of cut edges is $\frac m2$, and some split achieves at least this.
\end{proof}

\section{Alterations}

A random object may have a few defects; if they are few enough, we remove them.

\begin{proposition}\label{prop:alteration}
\begin{enumerate}[label=(\alph*)]
\item For every $n$, $R(k)>n-\binom nk2^{1-\binom k2}$.
\item A graph with $n$ vertices and average degree $d\ge1$ has an independent set (a set of vertices with no edge inside) of size at
least $\frac n{2d}$.
\end{enumerate}
\end{proposition}
\begin{proof}
(a) Colour $K_n$ at random and delete one vertex from each monochromatic $K_k$. What remains has no monochromatic $K_k$, and its expected
number of vertices is at least $n-\binom nk2^{1-\binom k2}$; some outcome does at least as well, and the number of remaining vertices is an integer.

(b) Keep each vertex with probability $p$, independently, and then delete one endpoint of every edge with both ends kept. The rest is
independent. The graph has $\frac{nd}2$ edges, so the expected number of vertices left is at least $pn-p^2\frac{nd}2$; with $p=\frac1d$ this is $\frac n{2d}$.
\end{proof}

For $k=10$, part (a) with $n=127$ gives $R(10)>115$, against $100$ from Theorem~\ref{thm:erdosramsey}; asymptotically, the alteration gains a factor $\sqrt2$.
Part (b) is a weak form of Tur\'an's theorem; Exercise~\ref{ex:carowei} gives a sharper one with a cleverer random choice.

\section{The second moment method}

The first moment shows that $X=0$ is possible. To show that $X>0$ is likely, we need the variance.

\begin{lemma}[second moment]\label{lem:secondmoment}
If $X$ is a random variable with $\mathbb E[X^2]<\infty$ and $\mathbb E[X]>0$, then $P(X=0)\le\frac{\operatorname{Var}(X)}{\mathbb E[X]^2}$.
\end{lemma}
\begin{proof}
If $X=0$ then $|X-\mathbb E X|\ge\mathbb E X$; apply Chebyshev's inequality.
\end{proof}

The \emph{random graph} $G(n,p)$ has vertex set $\{1,\dots,n\}$, and each of the $\binom n2$ possible edges is present independently with probability $p$.
As $p$ grows from $0$ to $1$, the graph acquires properties one after another, and many of them appear suddenly. We
say that an event holds \emph{with high probability} (w.h.p.) if its probability tends to $1$ as $n\to\infty$; here $p$ may depend on $n$.

\begin{theorem}[threshold for triangles]\label{thm:triangles}
Let $T$ be the number of triangles in $G(n,p)$. If $np\to0$, then w.h.p.\ $T=0$. If $np\to\infty$, then w.h.p.\ $T>0$.
\end{theorem}
\begin{proof}
For each set $A$ of three vertices let $I_A$ be the indicator that its three edges are present, so $T=\sum_AI_A$ and $\mathbb E[T]=\binom n3p^3\le\frac{(np)^3}6$. If $np\to0$,
then $P(T>0)\le\mathbb E[T]\to0$ by Markov's inequality.

Now let $np\to\infty$, so $\mathbb E[T]\sim\frac{(np)^3}6\to\infty$. Expand $\operatorname{Var}(T)=\sum_{A,B}\operatorname{Cov}(I_A,I_B)$. If $A$ and $B$ share at most one vertex, they have no edge in
common, so $I_A$ and $I_B$ are independent and the covariance is $0$. If $A=B$, it is at most $\mathbb E[I_A]=p^3$. If $A\ne B$ share two vertices, then
$A\cup B$ has four vertices and five edges, and the covariance is at most $\mathbb E[I_AI_B]=p^5$. Each set of four vertices contains four triangles, hence
$12$ ordered pairs of distinct ones, all sharing an edge. So
\[
\operatorname{Var}(T)\le\mathbb E[T]+12\binom n4p^5,
\]
and by the lemma
\[
\begin{aligned}P(T=0)&\le\frac1{\mathbb E[T]}+\frac{12\binom n4p^5}{\binom n3^2p^6}\eqbr\le\frac1{\mathbb E[T]}+\frac{18+o(1)}{n^2p}\to0,\end{aligned}
\]
since $\binom n4\big/\binom n3^2\sim\frac{36}{24n^2}$ and $n^2p\ge np\to\infty$.
\end{proof}

The function $\frac1n$ is a \emph{threshold} for containing a triangle. At the threshold itself, when $np=c$, the number of triangles converges in law to a
Poisson variable of mean $\frac{c^3}6$ (by the method of factorial moments of Chapter~2; see Alon and Spencer, Chapter~8), so $P(T>0)\to1-e^{-c^3/6}$. Every property preserved by adding edges
has a threshold (Bollob\'as and Thomason, 1987), and Exercise~\ref{ex:connectivity} finds the one for connectivity, $\frac{\ln n}n$.

\section{The local lemma}

If $A_1,\dots,A_n$ are bad events, the union bound shows that none occurs with positive probability when $\sum P(A_i)<1$. If they are
independent, it suffices that each $P(A_i)<1$. The local lemma interpolates: it suffices that each event is rare compared with the
number of events it depends on.

An event $A$ is \emph{mutually independent} of a family of events $\{A_j:j\in J\}$ if $P(A\cap C)=P(A)P(C)$ for every $C$ that is an intersection of events
$A_j$ or complements $A_j^c$, $j\in J$. Write $\bar A=A^c$.

\begin{theorem}[Lov\'asz local lemma, symmetric form]\label{thm:lll}
Let $A_1,\dots,A_n$ be events with $P(A_i)\le p$ for all $i$. Suppose that each $A_i$ is mutually independent of all but at most $d$ of the other
events, and that $ep(d+1)\le1$. Then
\[
P\Bigl(\bigcap_{i=1}^n\bar A_i\Bigr)\ge\Bigl(1-\frac1{d+1}\Bigr)^n>0.
\]
\end{theorem}
\begin{proof}
For each $i$, let $D_i$ be a set of at most $d$ indices such that $A_i$ is mutually independent of $\{A_j:j\notin D_i\cup\{i\}\}$. For $S\subseteq\{1,\dots,n\}$ write
$B_S=\bigcap_{j\in S}\bar A_j$ (with $B_\emptyset=\Omega$). We prove by induction on $|S|$ that $P(B_S)>0$ and that
\[
P(A_i\mid B_S)\le\frac1{d+1}\qquad\text{for every }i\notin S.\tag{$*$}
\]
For $S=\emptyset$, $P(A_i)\le p\le\frac1{e(d+1)}$. Let $|S|\ge1$ and suppose the claims hold for smaller sets. First, removing the elements of $S$ one at a
time, $P(B_S)=\prod_tP(\bar A_{s_t}\mid B_{\{s_1,\dots,s_{t-1}\}})\ge(1-\frac1{d+1})^{|S|}>0$ by $(*)$ for the smaller sets. Now fix $i\notin S$ and split $S$ into $S_1=S\cap D_i$ and $S_2=S\setminus D_i$.
If $S_1=\emptyset$, then $P(A_i\mid B_S)=P(A_i)\le p$ by mutual independence. Otherwise
\[
P(A_i\mid B_S)=\frac{P(A_i\cap B_{S_1}\mid B_{S_2})}{P(B_{S_1}\mid B_{S_2})}.
\]
The numerator is at most $P(A_i\mid B_{S_2})=P(A_i)\le p$, by mutual independence. For the denominator, list $S_1=\{j_1,\dots,j_r\}$ with $r\le d$; by $(*)$ applied to the
sets $S_2\cup\{j_1,\dots,j_{t-1}\}$, which are smaller than $S$,
\[
\begin{aligned}P(B_{S_1}\mid B_{S_2})&=\prod_{t=1}^r\Bigl(1-P\bigl(A_{j_t}\mid B_{S_2\cup\{j_1,\dots,j_{t-1}\}}\bigr)\Bigr)\eqbr\ge\Bigl(1-\frac1{d+1}\Bigr)^d\ge\frac1e,\end{aligned}
\]
since $(1+\frac1d)^d\le e$. So $P(A_i\mid B_S)\le ep\le\frac1{d+1}$. Finally, $P(B_{\{1,\dots,n\}})\ge(1-\frac1{d+1})^n$ as above.
\end{proof}

A \emph{$k$-CNF formula} is a conjunction of clauses, each a disjunction of $k$ literals (variables or their negations) on $k$ distinct
variables, such as $(x_1\vee\bar x_2\vee x_4)\wedge(\bar x_1\vee x_3\vee\bar x_4)$. Deciding whether such a formula can be satisfied is the archetypal NP-complete problem for
$k\ge3$.

\begin{corollary}\label{cor:ksat}
If each clause of a $k$-CNF formula shares a variable with at most $d$ other clauses, and $e(d+1)\le2^k$, then the formula is satisfiable.
\end{corollary}
\begin{proof}
Give the variables independent uniform values in $\{\text{true},\text{false}\}$, and let $A_i$ be the event that clause $i$ is false: all its $k$ literals are false,
probability $2^{-k}$. The event $A_i$ is determined by the variables of clause $i$, so it is mutually independent of the events of the clauses
that share no variable with it (Proposition~\ref{prop:indepfun}, functions of disjoint groups of independent variables). Theorem~\ref{thm:lll} with $p=2^{-k}$ applies.
\end{proof}

The union bound would require fewer than $2^k$ clauses in total; the local lemma allows any number, as long as each clause
meets fewer than $\frac{2^k}e$ others. For a long time the lemma gave no way to find the satisfying assignment, which is exponentially rare.
In 2009 Robin Moser found an astonishingly simple algorithm: start from a random assignment, and as long as some clause
is false, resample its variables at random. Moser and Tardos (2010) proved that under the hypothesis of the local
lemma the expected number of resamplings is at most $\frac n d$, where $n$ is the number of clauses. A numerical calculation runs it.

\section{The phase transition of random graphs}

Now let $p=\frac cn$ with a constant $c>0$: each vertex has about $c$ neighbours on average. Erd\H os and R\'enyi discovered in 1960 that
the structure of $G(n,\frac cn)$ changes abruptly at $c=1$ (Figure~\ref{fig:giant}). The key is to explore the component $C(v)$ of a vertex $v$ step by step.

\emph{Exploration.} Vertices are \emph{active}, \emph{explored} or \emph{neutral}. Initially $v$ is active and all others are neutral. At each step $i=1,2,\ldots$,
choose an active vertex $w$, reveal its edges to the neutral vertices, make the $X_i$ neutral neighbours found active, and mark
$w$ explored. The number of active vertices after step $i$ is $A_i=1+X_1+\dots+X_i-i$, the exploration stops at the first $i$ with $A_i=0$, and then $|C(v)|=i$.
Each edge is revealed at most once, so given the past, $X_i$ is binomial$(N_i,p)$, where $N_i\le n-1$ is the number of neutral vertices. Just as for a
Galton--Watson process with Poisson$(c)$ offspring (Chapter~3), the exploration should die out if $c<1$ and may survive if $c>1$.

\begin{theorem}[small components]\label{thm:subcritical}
Let $c<1$ and $I=c-1-\ln c>0$. For every $a>\frac1I$, w.h.p.\ every component of $G(n,\frac cn)$ has at most $a\ln n$ vertices.
\end{theorem}
\begin{proof}
\emph{Coupling.} At step $i$, add to the $N_i$ trials of the exploration $n-N_i$ independent ``ghost'' trials of probability $p$, and let $\tilde X_i\ge X_i$ count the
successes among all $n$ trials (after the exploration has stopped, let $\tilde X_i$ count $n$ fresh trials). Given the past, $\tilde X_i$ is binomial$(n,p)$, so the $\tilde X_i$
are independent binomial$(n,p)$ variables, and $\tilde X_1+\dots+\tilde X_k$ is binomial$(nk,p)$, with mean $ck$.

\emph{One vertex.} If $|C(v)|>k$, the exploration is still running after $k$ steps, so $A_k\ge1$, that is $X_1+\dots+X_k\ge k$, and so $\tilde X_1+\dots+\tilde X_k\ge k$. By the
Chernoff bound (Theorem~\ref{thm:chernoff}) with $\mu=ck$ and $1+\delta=\frac1c$,
\[
\begin{aligned}P(|C(v)|>k)&\le\Bigl(\frac{e^\delta}{(1+\delta)^{1+\delta}}\Bigr)^{ck}\eqbr=e^{-k(c-1-\ln c)}=e^{-kI}.\end{aligned}
\]
\emph{All vertices.} With $k=\lfloor a\ln n\rfloor$, the union bound over the $n$ vertices gives
\[
P\Bigl(\max_v|C(v)|>k\Bigr)\le ne^{-kI}\le e^In^{1-aI}\to0.
\]
\end{proof}

For $c>1$ the exploration from a vertex survives with probability close to the survival probability $\rho$ of the Poisson$(c)$ branching
process of Exercise~\ref{ex:poissongw}, and the surviving explorations all end up in a single component.

\begin{remark}
For $c>1$, w.h.p.\ $G(n,\frac cn)$ has a unique \emph{giant component}, with $(\rho+o(1))n$ vertices, where $\rho\in\,]0,1[$ solves $1-\rho=e^{-c\rho}$, while all other
components have $O(\ln n)$ vertices (Erd\H os and R\'enyi, 1960). Exercise~\ref{ex:giant} proves the existence of a component of linear size by a
depth-first search argument of Krivelevich and Sudakov; the uniqueness and the exact constant $\rho$ need finer couplings, found in
the books of Janson, \L uczak and Ruci\'nski, of Frieze and Karo\'nski, and of van der Hofstad. A numerical calculation watches $\rho$ appear. At $c=1$
itself the largest component has about $n^{2/3}$ vertices, a critical window studied by Aldous with Brownian motion.
\end{remark}

\begin{figure}[ht]
\centering
\fig{ch12-giant}
\caption{The largest (blue) and second largest (orange) components of a random graph on $n=20\,000$ vertices with $\frac{cn}2$ random edges (essentially
$G(n,\frac cn)$), as fractions of $n$, against $c$; edges are added one by one, so that all values of $c$ come from one growing graph.
The black curve is the survival probability $\rho(c)$ of a Poisson$(c)$ branching process.}
\label{fig:giant}
\end{figure}

\begin{beauty}
A random graph knows about branching processes. Seen from one vertex, $G(n,\frac cn)$ looks like a Poisson$(c)$ family tree: the number
of neighbours is close to Poisson$(c)$, and so is the number of new neighbours of each of them, as long as the exploration is
small compared with $n$. The extinction probability of Galton's family names (Chapter~3) becomes the fraction of
vertices outside the giant component.
\end{beauty}

\section{From existence to algorithms}

A probabilistic proof often hides an algorithm. Choose the random bits one at a time, each time keeping the
conditional expectation of the quantity of interest from decreasing. When these conditional expectations can be
computed, this is the \emph{method of conditional expectations}.

\begin{theorem}\label{thm:maxsat}
For a $k$-CNF formula with $m$ clauses, there is an assignment satisfying at least $(1-2^{-k})m$ clauses, and the following deterministic
algorithm finds one: for $i=1,\dots,n$, give $x_i$ the value that maximises the expected number of satisfied clauses when $x_1,\dots,x_i$ are fixed
and the remaining variables are uniform.
\end{theorem}
\begin{proof}
Let $W$ be the number of satisfied clauses under a uniform random assignment and $W_i(a_1,\dots,a_i)$ its conditional expectation given
$x_1=a_1,\dots,x_i=a_i$, that is, $\mathbb E[W\mid x_1,\dots,x_i]$ evaluated at these values (Chapter~10, for a finite partition). A clause contributes $1$ if one of its fixed
literals is true, and otherwise $1-2^{-u}$, where $u$ is its number of unfixed literals ($0$ if $u=0$): each $W_i$ is computable in time proportional
to the size of the formula. By the tower property, writing $a=(a_1,\dots,a_{i-1})$,
\[
\begin{aligned}W_{i-1}(a)&=\tfrac12W_i(a,\text{true})\eqbr+\tfrac12W_i(a,\text{false}),\end{aligned}
\]
so the better of the two values satisfies $W_i\ge W_{i-1}$. Starting from $W_0=\mathbb E[W]=(1-2^{-k})m$, the algorithm ends with a full assignment satisfying
$W_n\ge(1-2^{-k})m$ clauses.
\end{proof}

For $k=3$ this gives $\frac78$ of the clauses. H\aa stad proved in 2001 that, unless P $=$ NP, no polynomial algorithm can guarantee $\frac78+\varepsilon$ for
satisfiable $3$-CNF formulas: the trivial random assignment, derandomised, is optimal. For MAX-CUT, the same method
places the vertices one by one on the side that cuts more edges to the already placed ones, and finds a cut of at least
$\frac m2$ edges (Proposition~\ref{prop:maxcut}).

\section{Exercises}

\begin{exercise}[$\star$]\label{ex:carowei}
\emph{The Caro--Wei bound.} Order the vertices of a graph uniformly at random, and keep the vertices that come before all their
neighbours. Show that this is an independent set, and deduce that the graph has an independent set of size at least
$\sum_v\frac1{d(v)+1}$, where $d(v)$ is the degree of $v$. Compare with Proposition~\ref{prop:alteration}(b).
\end{exercise}

\begin{exercise}[$\star$]\label{ex:tournament}
\emph{Tournaments.} In a round-robin tournament of $n$ players every pair plays once, with no draws. Show that if
$\binom nk(1-2^{-k})^{n-k}<1$, there is a tournament in which every set of $k$ players is beaten by some other player (who beats all $k$ of
them). Deduce that such tournaments exist for $k=2$ with $n=21$ players and for $k=3$ with $n=91$.
\end{exercise}

\begin{exercise}[$\star\star$]\label{ex:sumfree}
\emph{Sum-free sets (Erd\H os, 1965).} A set of integers is \emph{sum-free} if it contains no $x$, $y$, $z$ with $x+y=z$. Show that every set $A$ of $n$
non-zero integers contains a sum-free subset of more than $\frac n3$ elements. (Take a prime $p=3q+2$ larger than all $|a|$, $a\in A$; for a uniform
$t\in\{1,\dots,p-1\}$, keep the $a\in A$ such that $ta\bmod p$ lies in the middle third $\{q+1,\dots,2q+1\}$.)
\end{exercise}

\begin{exercise}[$\star\star$]\label{ex:crossing}
\emph{The crossing lemma.} The crossing number $\operatorname{cr}(G)$ of a graph is the least number of edge crossings in a drawing of it in the plane.
Using Euler's formula for planar graphs (a planar graph with $n\ge3$ vertices has at most $3n-6$ edges), show that $\operatorname{cr}(G)\ge m-3n$ for
every graph with $n$ vertices and $m$ edges. Then apply this to a random induced subgraph to show that $\operatorname{cr}(G)\ge\frac{m^3}{64n^2}$ when $m\ge4n$.
\end{exercise}

\begin{exercise}[$\star$]\label{ex:propertyb}
\emph{Colouring hypergraphs.} A family of $k$-element sets is \emph{2-colourable} if the elements can be coloured red and blue so that no set
of the family is monochromatic. Show that a family of fewer than $2^{k-1}$ sets is 2-colourable, and that so is any family in which each
set meets at most $d$ others, if $e(d+1)\le2^{k-1}$.
\end{exercise}

\begin{exercise}[$\star\star\star$]\label{ex:connectivity}
\emph{The connectivity threshold.} Let $p=\frac{\ln n+\omega}n$, where $\omega=\omega(n)$.
(a) Let $Y$ be the number of isolated vertices. Show that $\mathbb E[Y]\sim e^{-\omega}$ if $\omega=o(\ln n)$, and that if $\omega\to-\infty$, then w.h.p.\ $Y>0$, so the graph is not connected.
(b) Suppose $\omega\to+\infty$ with $\omega\le\ln\ln n$. Show that the expected number of sets $S$ of $k$ vertices, $1\le k\le\frac n2$, with no edge between $S$ and its complement
and spanning a connected subgraph, tends to $0$. Conclude that w.h.p.\ the graph is connected. (A connected graph on $k$ vertices contains a
spanning tree, and there are $k^{k-2}$ trees on $k$ labelled vertices, by Cayley's formula.) Explain why the condition $\omega\le\ln\ln n$ may be dropped.
\end{exercise}

\begin{exercise}[$\star\star\star$]\label{ex:giant}
\emph{A giant component, by depth-first search.} Let $p=\frac{1+\varepsilon}n$ with $0<\varepsilon\le\frac1{10}$. Run a depth-first search on $G(n,p)$, maintaining three sets: $S$
(finished vertices), $U$ (a stack) and $T$ (unvisited vertices). If $U$ is empty, move some vertex of $T$ to $U$. Otherwise, let $v$ be the top of
$U$, and query the pairs $\{v,w\}$, $w\in T$, one at a time: at the first edge found, push $w$ onto $U$; if there is none, move $v$ to $S$. The answers to
the queries are independent Bernoulli$(p)$ variables.
(a) Show that the vertices of $U$ always lie on a path of the graph, that every pair with one vertex in $S$ and one in $T$ has been
queried, and that $|S|+|U|$ is at least the number of positive answers so far.
(b) Let $N=\lfloor\frac{\varepsilon n^2}2\rfloor$. Show that w.h.p.\ the first $N$ answers contain at least $\frac{\varepsilon(1+\varepsilon)n}2-n^{2/3}$ positive ones.
(c) Show that after $N$ queries, $|S|<\frac n3$. (Otherwise consider the moment when $|S|=\frac n3$.)
(d) Deduce that after $N$ queries, w.h.p.\ $|U|\ge\frac{\varepsilon^2n}5$. Conclude that w.h.p.\ $G(n,p)$ has a component, and even a path, with at least $\frac{\varepsilon^2n}5$ vertices,
and that the same holds for $G(n,\frac cn)$ for every $c>1$.
\end{exercise}

\section{Hints}

\begin{hint}{ex:carowei}
$v$ is kept with probability $\frac1{d(v)+1}$: it must come first among itself and its $d(v)$ neighbours.
\end{hint}
\begin{hint}{ex:tournament}
Orient each game by a fair coin. A fixed set of $k$ players is beaten by a fixed outsider with probability $2^{-k}$.
\end{hint}
\begin{hint}{ex:sumfree}
The middle third is sum-free modulo $p$, and for each $a\ne0$, $ta\bmod p$ is uniform on $\{1,\dots,p-1\}$.
\end{hint}
\begin{hint}{ex:crossing}
Remove one edge per crossing. For the random subgraph, keep each vertex with probability $q=\frac{4n}m$ and take expectations of
$\operatorname{cr}\ge m-3n$ for the subgraph with the induced drawing.
\end{hint}
\begin{hint}{ex:propertyb}
A given set is monochromatic with probability $2^{1-k}$.
\end{hint}
\begin{hint}{ex:connectivity}
(a) $\mathbb E[Y]=n(1-p)^{n-1}$ and $\mathbb E[Y(Y-1)]=n(n-1)(1-p)^{2n-3}$. (b) Bound the expectation by $\sum_k\binom nkk^{k-2}p^{k-1}(1-p)^{k(n-k)}$ and use $1-p\le e^{-p}$. For the last
question, connectivity is preserved by adding edges.
\end{hint}
\begin{hint}{ex:giant}
(b) Chernoff. (c) At that moment $|U|\le1+\varepsilon n$ and all pairs between $S$ and $T$ have been queried. (d) If $|U|<\frac{\varepsilon^2n}5$, then $|S|$ is large, and
$|S|\cdot|T|$ exceeds $N$.
\end{hint}

\section{Solutions}

\begin{solution}{ex:carowei}
\idea A random order resolves every edge in favour of one endpoint, and a vertex wins against all its neighbours with
probability one over the size of its neighbourhood.
\writeup If two kept vertices were adjacent, each would come before the other. Vertex $v$ is kept when it is first among the $d(v)+1$
vertices of its closed neighbourhood, which happens with probability $\frac1{d(v)+1}$ by symmetry. By linearity the expected size of the kept set is
$\sum_v\frac1{d(v)+1}$, and some order achieves at least this. By convexity of $x\mapsto\frac1{x+1}$ (Jensen, Theorem~\ref{thm:jensenprob} for the uniform law on the vertices), this is at
least $\frac n{d+1}$, where $d$ is the average degree: about twice the bound $\frac n{2d}$ of Proposition~\ref{prop:alteration}(b). For a disjoint union of cliques $K_{d+1}$ it is exact.
\end{solution}

\begin{solution}{ex:tournament}
\idea Count the bad events: a set of $k$ players that nobody beats entirely. Each outsider fails independently.
\writeup Orient each game by an independent fair coin. For a set $K$ of $k$ players and an outsider $w$, $w$ beats all of $K$ with probability $2^{-k}$, and
these events are independent for the $n-k$ outsiders, which involve disjoint sets of games. So $K$ is unbeaten with probability $(1-2^{-k})^{n-k}$,
and the expected number of unbeaten sets is $\binom nk(1-2^{-k})^{n-k}<1$: some tournament has none. Numerically, $\binom{21}2(\frac34)^{19}\approx0.89$ and $\binom{91}3(\frac78)^{88}\approx0.96$. (The smallest
tournament for $k=2$ has $7$ players: let $i$ beat $j$ when $j-i$ is a non-zero square modulo $7$.)
\end{solution}

\begin{solution}{ex:sumfree}
\idea Multiplying by a random unit spreads each element uniformly over the non-zero residues, and the middle third is a
sum-free target of relative size just above one third.
\writeup Let $M=\{q+1,\dots,2q+1\}$, which has $q+1$ elements. If $x,y\in M$ then $2q+2\le x+y\le4q+2$, and modulo $p=3q+2$ the sum lies in $\{2q+2,\dots,3q+1\}\cup\{0,\dots,q\}$, outside $M$:
$M$ is sum-free modulo $p$. For $a\in A$, $a\not\equiv0\pmod p$ since $p>|a|$, so $t\mapsto ta\bmod p$ is a bijection of $\{1,\dots,p-1\}$ ($\Z/p\Z$ is a field), and $ta\bmod p$ is uniform: it lies in
$M$ with probability $\frac{q+1}{3q+1}>\frac13$. So the expected number of kept elements exceeds $\frac n3$, and some $t$ keeps more than $\frac n3$. The kept set $A_t$ is
sum-free: if $x+y=z$ in $A_t$, then $tx+ty\equiv tz\pmod p$ with all three residues in $M$, which is impossible. (Infinitely many primes $p\equiv2\pmod3$ exist, by
the argument of Euclid applied to $3N!-1$.)
\end{solution}

\begin{solution}{ex:crossing}
\idea A drawing with few crossings would give random subgraphs with fewer crossings than their many edges force.
\writeup Take a drawing with $\operatorname{cr}(G)$ crossings and delete one edge from each crossing: the remaining graph is drawn without crossings, so it is
planar and has at most $3n$ edges; hence $m-\operatorname{cr}(G)\le3n$ (also trivially true for $n\le2$). Now keep each vertex with probability $q$, independently, and consider the
induced subgraph $G'$ with the induced drawing, with $n'$ vertices, $m'$ edges and $c'$ crossings in that drawing. A crossing involves two edges with four
distinct endpoints (in an optimal drawing, edges with a common endpoint do not cross), so $\mathbb E[c']=q^4\operatorname{cr}(G)$, while $\mathbb E[m']=q^2m$ and $\mathbb E[n']=qn$. As $c'\ge\operatorname{cr}(G')\ge m'-3n'$ for
every outcome, $q^4\operatorname{cr}(G)\ge q^2m-3qn$. With $q=\frac{4n}m\le1$:
\[
\operatorname{cr}(G)\ge\frac m{q^2}-\frac{3n}{q^3}=\frac{m^3}{16n^2}-\frac{3m^3}{64n^2}=\frac{m^3}{64n^2}.
\]
This short proof, from the book of Aigner and Ziegler, has applications from incidence geometry (the Szemer\'edi--Trotter
theorem) to sum-product estimates.
\end{solution}

\begin{solution}{ex:propertyb}
\idea Exactly as for Ramsey numbers and for $k$-SAT: a random colouring, then the union bound or the local lemma.
\writeup Colour each element red or blue at random. A given set is monochromatic with probability $2\cdot2^{-k}=2^{1-k}$. With fewer than $2^{k-1}$ sets, the expected
number of monochromatic sets is less than $1$. In the second case, the event that a set is monochromatic is mutually independent of the
events for the sets disjoint from it (Proposition~\ref{prop:indepfun}), and Theorem~\ref{thm:lll} with $p=2^{1-k}$ applies.
\end{solution}

\begin{solution}{ex:connectivity}
\idea Below the threshold the graph still has isolated vertices; above it, the only way to be disconnected would be a
small piece cut off, and these are too unlikely.
\writeup (a) $\mathbb E[Y]=n(1-p)^{n-1}$. Since $\ln(1-p)=-p+O(p^2)$ and $np^2\to0$, $(1-p)^{n-1}=e^{-np+o(1)}=\frac{e^{-\omega}}ne^{o(1)}$, so $\mathbb E[Y]\sim e^{-\omega}$ (when $p\to0$, which holds if $\omega=o(\ln n)$).
Two given vertices are both isolated with probability $(1-p)^{2n-3}$ (they avoid $2n-3$ edges), so $\mathbb E[Y^2]=\mathbb E[Y]+n(n-1)(1-p)^{2n-3}$ and
\[
\begin{aligned}\frac{\operatorname{Var}(Y)}{\mathbb E[Y]^2}&\le\frac1{\mathbb E[Y]}+\frac{(1-p)^{-1}n(n-1)}{n^2}-1\eqbr\le\frac1{\mathbb E[Y]}+\frac p{1-p}.\end{aligned}
\]
If $\omega\to-\infty$ (with $\omega\ge-\ln n$, so that $p\ge0$), then $p\to0$ and, using $\ln(1-p)\ge-p-p^2$ for $p\le\frac12$,
$\mathbb E[Y]\ge n(1-p)^n\ge e^{\ln n-np-np^2}=e^{-\omega-np^2}\to\infty$, so $P(Y=0)\to0$ by Lemma~\ref{lem:secondmoment}: w.h.p.\ there is an isolated vertex, and the graph is disconnected.

(b) If the graph is disconnected, some component $S$ has $1\le|S|\le\frac n2$ vertices: it is connected, so contains one of the $k^{k-2}$ spanning trees on $S$, whose
$k-1$ edges are present, and the $k(n-k)$ pairs between $S$ and its complement are absent. The expected number of such pairs (set, tree) is
\[
E=\sum_{k=1}^{n/2}\binom nkk^{k-2}p^{k-1}(1-p)^{k(n-k)}.
\]
Using $\binom nk\le\frac{n^k}{k!}\le\bigl(\frac{en}k\bigr)^k$, $k^{k-2}\le k^k$, $(1-p)^{k(n-k)}\le e^{-pk(n-k)}$ the $k$-th term is at most $\frac1p\bigl(enp\,e^{-p(n-k)}\bigr)^k$, that is
\[
\frac1p\Bigl(e(\ln n+\omega)\,e^{-(\ln n+\omega)(1-\frac kn)}\Bigr)^k.
\]
For $k\le\frac n{\ln n}$, $(\ln n+\omega)\frac kn\le2$, so the base is at most $\frac{e^3(\ln n+\omega)e^{-\omega}}n$ and the sum over these $k$ is at most $\frac1p\cdot2\frac{e^3(\ln n+\omega)e^{-\omega}}n\le2e^3e^{-\omega}\to0$
(a geometric series with ratio $\le\frac12$ for large $n$, and $\frac1p\cdot\frac{\ln n+\omega}n=1$). For $\frac n{\ln n}<k\le\frac n2$, the base is at most $e(\ln n+\omega)e^{-(\ln n+\omega)/2}\le\frac{e\cdot2\ln n}{\sqrt n}$, which is less
than $n^{-1/3}$ for large $n$, so these terms add up to at most $\frac n{p}\,n^{-\frac{n}{3\ln n}}\to0$. So $E\to0$, and by Markov's inequality the graph is connected w.h.p.
Finally, if $\omega\to\infty$ faster, replace $\omega$ by $\min(\omega,\ln\ln n)$: the graph $G(n,p)$ can be obtained from $G(n,p')$ with $p'\le p$ by adding
independent edges (keep each absent edge with probability $\frac{p-p'}{1-p'}$), and adding edges preserves connectivity. Together with (a), $\frac{\ln n}n$ is a
sharp threshold: with $\omega=c$ constant, $P(\text{connected})\to e^{-e^{-c}}$, the probability that a Poisson$(e^{-c})$ number of isolated vertices is zero (Erd\H os and R\'enyi, 1959).
\end{solution}

\begin{solution}{ex:giant}
\idea Depth-first search keeps its current branch on a stack, which is a path. If the stack stayed short, many
vertices would be finished early, and all the pairs between them and the unvisited vertices would have been queried
with no success: too many queries for the time available.
\writeup (a) A vertex is pushed onto $U$ when it is adjacent to the current top, so $U$ is always a path, bottom to top, and lies in one component.
When $v$ moves to $S$, all pairs $\{v,w\}$ with $w$ in $T$ at that moment have been queried, negatively; $T$ only shrinks, so all pairs between $S$ and $T$ have
been queried. Each positive answer moves a new vertex from $T$ to $U$, and vertices never return to $T$, so $|S|+|U|$ is at least the number of positive
answers. Moreover $|U|$ exceeds by at most $1$ the number of positive answers since $U$ was last empty.

(b) The number $P_N$ of positive answers among the first $N$ is binomial$(N,p)$, with mean $Np=\frac{\varepsilon(1+\varepsilon)n}2+O(1)$. By Chernoff (Theorem~\ref{thm:chernoff}) with $\delta=n^{2/3}/(Np)$,
$P(P_N\le Np-n^{2/3})\le e^{-n^{4/3}/(2Np)}\to0$, since $Np$ is of order $n$.

(c) The search makes at most $\binom n2>N$ queries before it ends, so it is still running at time $N$. Suppose $|S|\ge\frac n3$ at time $N$; consider the first moment
when $|S|=\lceil\frac n3\rceil$. Then $|U|\le1+P_N$ by (a), and w.h.p.\ $P_N\le2Np\le1.1\varepsilon n$ (Chernoff again), so
$|T|\ge n-\frac n3-1-1.1\varepsilon n\ge\frac n2$ for large $n$ (as $\varepsilon\le\frac1{10}$). All $|S|\,|T|\ge\frac{n^2}6$ pairs between $S$ and $T$ have been queried, but $\frac{n^2}6>\frac{\varepsilon n^2}2\ge N$: a contradiction.

(d) Suppose $|U|<\frac{\varepsilon^2n}5$ at time $N$. By (a) and (b), $|S|\ge P_N-|U|\ge\frac{\varepsilon n}2+\frac{3\varepsilon^2n}{10}-n^{2/3}$, and by (c) $|S|<\frac n3$. Then
$|T|=n-|S|-|U|$, and since $s\mapsto s(n-s-|U|)$ increases for $s\le\frac{n-|U|}2$,
\[
\begin{gathered}|S|\,|T|\ge s_0\Bigl(n-s_0-\frac{\varepsilon^2n}5\Bigr),\qq s_0=\frac{\varepsilon n}2+\frac{3\varepsilon^2n}{10}-n^{2/3},\end{gathered}
\]
which equals $n^2\bigl(\frac\varepsilon2+\frac{\varepsilon^2}{20}-\frac{2\varepsilon^3}5-\frac{3\varepsilon^4}{20}\bigr)+O(n^{5/3})$. For $\varepsilon\le\frac1{10}$, $\frac{\varepsilon^2}{20}>\frac{2\varepsilon^3}5+\frac{3\varepsilon^4}{20}$, so for large $n$ this
exceeds $\frac{\varepsilon n^2}2\ge N$, although all these pairs have been queried: a contradiction. So w.h.p.\ $|U|\ge\frac{\varepsilon^2n}5$, and $U$ is a path in one component. For
$G(n,\frac cn)$ with $c>1$, take $\varepsilon=\min(c-1,\frac1{10})$: $G(n,\frac cn)$ contains a copy of $G(n,\frac{1+\varepsilon}n)$ (add independent edges as in Exercise~\ref{ex:connectivity}), and adding edges
does not shrink components.
\end{solution}

\section*{Sources and further reading}

\begin{itemize}
\item N.~Alon and J.~Spencer, \emph{The Probabilistic Method}, 4th ed., Wiley, 2016. The reference for the whole chapter.
\item M.~Aigner and G.~Ziegler, \emph{Proofs from THE BOOK}, 6th ed., Springer, 2018. Chapters on the probabilistic method and
on the crossing lemma.
\item P.~Erd\H os, \emph{Some remarks on the theory of graphs}, Bulletin of the American Mathematical Society 53 (1947),
292--294.
\item P.~Erd\H os and A.~R\'enyi, \emph{On the evolution of random graphs}, Publications of the Mathematical Institute of the
Hungarian Academy of Sciences 5 (1960), 17--61.
\item L.~Lov\'asz and P.~Erd\H os, \emph{Problems and results on 3-chromatic hypergraphs and some related questions}, in
Infinite and Finite Sets, Colloquia Mathematica Societatis J\'anos Bolyai 10, North-Holland, 1975, 609--627.
\item R.~Moser and G.~Tardos, \emph{A constructive proof of the general Lov\'asz local lemma}, Journal of the ACM 57 (2010),
article~11.
\item M.~Krivelevich and B.~Sudakov, \emph{The phase transition in random graphs: a simple proof}, Random Structures and
Algorithms 43 (2013), 131--138.
\item S.~Janson, T.~\L uczak and A.~Ruci\'nski, \emph{Random Graphs}, Wiley, 2000; A.~Frieze and M.~Karo\'nski, \emph{Introduction to
Random Graphs}, Cambridge University Press, 2016; R.~van der Hofstad, \emph{Random Graphs and Complex Networks}, Vol.~1,
Cambridge University Press, 2017.
\item M.~Campos, S.~Griffiths, R.~Morris and J.~Sahasrabudhe, \emph{An exponential improvement for diagonal Ramsey},
preprint, 2023, arXiv:2303.09521.
\item J.~H\aa stad, \emph{Some optimal inapproximability results}, Journal of the ACM 48 (2001), 798--859.
\item M.~Mitzenmacher and E.~Upfal, \emph{Probability and Computing}, 2nd ed., Cambridge University Press, 2017, Chapter~6.
\end{itemize}
